%Usage guide: https://docs.google.com/document/d/1oDJKcFVK_yXD3Do6vfyTFKJuDwdwN4g6XFTM5QASadA/edit?usp=sharing

\documentclass[12pt]{article} % If you want 12 pt font instead of the standard 10 pt
%\documentclass{article}
\usepackage[utf8]{inputenc}


\usepackage{amsmath}
\usepackage{amssymb}
\usepackage[shortlabels]{enumitem}
\usepackage[margin=1 in]{geometry}

\pagenumbering{arabic}

\usepackage{mathhw}

\begin{document}
\note{L17 - Linear Regression}

Given a sample $(X_i,Y_i)$, i.i.d. for some unknown joint distribution $\statP$, then $\statP$ is described entirely by either the joint PDF $h(x,y)$, or the marginal density $h(x)$ AND the corresponding conditional density for $h(y|x)$.

Then the regression function $x \rightarrow f(x)$ is of the form:
\begin{equation*}
    E[Y|X = x] = \int y\ h(y|x)\ dy
\end{equation*}

\pinkbox{
The \textbf{theoretical linear regression} of $Y$ on $X$ is the line $a^* + b^*x$ with:
\begin{equation*}
    (a^*,b^*) = \arg \min E[(Y-a-bX)^2]
\end{equation*}

where setting partial derivatives to zero gives
\begin{align*}
    b^* &= \frac{\Cov{X}{Y} }{\Var(X)} \\
    a^* &= E[Y] - b^*E[X]
\end{align*}
clearly there will be some noise $\epsilon$ such that $Y - (a^*+b^*X) = \epsilon$. This noise has
\begin{itemize}[topsep=-1pt,itemsep=-1pt]
    \item $E[\epsilon] = 0$
    \item $\Cov{X}{\epsilon} = 0$
\end{itemize}


The \textbf{Least Squares Estimator} of $a^*,b^*$ is the minimizer of the sum of squared errors, and is given by:
\begin{align*}
    \hat{b} &= \frac{\overline{XY}- \overline{X}\overline{Y}}{\overline{X^2} - \overline{X}^2} \\
    \hat{a} &= \overline{Y} - \hat{b}\overline{X}
\end{align*}
}



\end{document}


\begin{comment}
Pink - Important ideas/defs/thsm
Green - Corollaries/lesser thms 
Blue - Smaller words
Yellow - Examples
Gray - Anything trivial

\begin{itemize}[topsep=-1pt,itemsep=-1pt]
    \item Hi
\end{itemize}


\begin{enumerate}[label=\textbf{(\alph*)}, topsep=0pt]
    \item Hello
\end{enumerate}


\begin{center}
   \includegraphics[width=0.6\textwidth]{L15-1.png} 
\end{center}
\end{comment} 